% Job posting: https://ca.linkedin.com/jobs/view/research-engineer-%E2%80%93-multimodal-ai-interaction-at-huawei-canada-4463665160
\documentclass[11pt]{article}
\usepackage[left=1in,top=1in,right=1in,bottom=1in]{geometry}
\usepackage[hidelinks]{hyperref}
\usepackage{setspace}
\usepackage[T1]{fontenc}
\usepackage{newtxtext,newtxmath}
\ifdefined\pdfgentounicode
  \input{glyphtounicode}
  \pdfgentounicode=1
\fi
\setstretch{1.1}
\pagenumbering{gobble}
\begin{document}
\pagestyle{empty}

\begin{flushright}
Nishal Stanislaus Silva, PhD \\
Toronto, ON \\
nishal.silva@hotmail.com \\
+1 613 200 2030
\end{flushright}

Dear Hiring Manager,

During my postdoctoral research at the University of Trento, I designed and led a real-time multimodal synchronization pipeline that fused streaming audio, BCI/EEG biosignals captured with g.tec hardware, and mixed-reality head and hand tracking across four simultaneous performers, developed in collaboration with SAE Institute Barcelona. I took the system beyond the lab: we co-designed and ran two live multisensory concerts built on it, coordinating stage lighting and smoke, ten Meta Quest 3 headsets, and haptic phones in real time, and evaluated the result through a structured user study of 20 audience members and 6 performers. That project is, at its core, multimodal AI and human-interaction research conducted at performance scale, and it is the piece of my work I believe maps most directly onto this role.

That project did not stand alone. As a visiting researcher at McGill University's IDMIL/CIRMMT lab, I fused IMU and audio streams through a hierarchical finite-state machine to build a self-powered smart electric guitar, reaching F1 0.78 on a 500-event gesture corpus. At the University of Visual and Performing Arts in Colombo, I ran a human-centred evaluation of musicians' perception of smart instruments, combining quantitative and qualitative methods. And throughout my PhD, I built the habit of validating real-time systems rigorously: frozen-protocol benchmark suites, latency/accuracy/CPU ablations, and open datasets, culminating in an audio pattern detector that runs in 14 ms on a Raspberry Pi 4 at F1 0.76, 74 times faster than the DTW baseline it replaced. Together, these give me a research process built specifically around fusing multiple live data streams and testing what they mean for the humans using them.

I have not worked with Huawei's specific consumer hardware or product ecosystem, and I want to be upfront about that. What I bring instead is a multimodal-fusion architecture and a human-interaction research methodology that transfer directly to a new hardware context, along with the speed to get there: I have built comparable pipelines from scratch across audio, IMU, BCI/EEG, and XR tracking more than once.

I am a Canadian Permanent Resident based in Toronto, eligible to work anywhere in Canada without sponsorship, which places me close to your Markham office. I am available immediately and would welcome the chance to discuss how this research background could contribute to your team.

Sincerely, \\
Nishal Stanislaus Silva, PhD

\end{document}
